

New emergent technologies use neurons in the area of computation, creating circuits that depend on these specialized cells to solve problems. these problems would be more difficult to implement or have higher computational cost. Different modern Real-Time Operative Systems (RTOS) and other tools can simplify the interaction between neurons; thus, some options are proposed, and with them, some tests for the performance and latency.

The goal of the interactions should be forming a neural hybrid circuit, with a biological part (live neuron) and an artificial one (model neuron), so it can be used to improve both technological and research fields. To achieve that, the time restriction that those interactions demand need to be respected. If not, the connection between the biological part and artificial part will not succeed. To comply with this restrictions, it is mandatory to use systems and tools specialized on this kind of problems, specialized on real-time.

In consequence, it is initially presented the formalization and definition, from the electrophysiology field, of the term \textit{hybrot} for its use in real-time, along with a small recapitulation of actual real-time tools. After this, the specific use of some of the tools and systems will be explored in depth, specifically one of the most used in the neuroscience community, \textit{RTXI}. This tool allows the execution of electrophysiology experiments in real-time. This program has been adapted during this work for its use in \textit{Preempt-RT}, the official real-time kernel for Linux. There are also presented a series of validations in form of latency tests and creation of biohybrid circuits to check its correctness in experimental neuroscience.

The latency and performance tests have been realized both in \textit{RTXI} and in \textit{cyclictest}, which is a program from the \textit{rt-test} suite of \textit{The Linux Foundation}. As for the biohybrid circuits, there have been done two types of validations. The first one is with an electronic neuron (implemented in an analog circuit) that presents the same time constraints as a live neuron. The second validation is done from an electrophisioloigcal register of a living neuron that acts as a stimulus to a neural model executing in the developed system.